\scanpage{3}
\mspage{\handwritten{6}}
\margin{\handwritten{Co Exp $\overline{\text{\uline{XIII}}}$}}

\uline{Proposition} 1.3. Avec les notations précédentes, \deleted{\unsure{en|on}} supp\handwritten{o}s\deleted{\unsure{ant|e}}\handwritten{ons} \deleted{encore} \deleted{\unsure{la|la en}} l'inégalité (2) valable\deleted{,}\inserted{et $G,V,W$ géométriquement irréductibles.} \deleted{c}\handwritten{C}onsidérons les conditions suivantes:

(i) $\handwritten{N}=\handwritten{M}_{\deleted{\unsure{a}}a}$, et $N$ est lisse sur $k$.

(ii) $\handwritten{N}=\handwritten{M}_{a}$, et $\dim X=\dim V$ (cf lemme 1.2.).

(iii) $\handwritten{N}=\handwritten{M}_{a}$, et $\handwritten{\psi}:X\handwritten{\to}V$ est génériquement quasi-fini.

(iii bis) $\handwritten{N}=\handwritten{M}_{a}$, et $\handwritten{\psi}:X\handwritten{\to}V$ est génériquement non ramffié.

(iii ter) $\handwritten{N}=\handwritten{M}_{a}$, et $\handwritten{\psi}:X\handwritten{\to}V$ est génériquement étale.

(iv) $\handwritten{\psi}:X\handwritten{\to}V$ est non ramffié en $(e^{\handwritten{\unsure{\circ|\bullet}}},a)$.

(iv bis) $\handwritten{\psi}:X\handwritten{\to}V$ est étale en $(e^{\handwritten{\circ}},a)$.

(v) $\handwritten{\psi}:X\handwritten{\to}V$ est lisse en $(e^{\handwritten{\circ}},a)$, et $N$ est lisse sur $k$, (2) une \handwritten{égalité.} \deleted{égalité}

(vi) $M_a$ et $N$ coïncident au voisinage de $e$.

Alors on a les implications
\[
\resizebox{.88\linewidth}{!}{$\displaystyle
\text{(\deleted{\unsure{i}}v)}\handwritten{\Leftrightarrow}\text{(iv bis)}\handwritten{\Longrightarrow}
\text{(iii \deleted{bis}ter)}\handwritten{\Longrightarrow}\text{(i)}\handwritten{\Leftrightarrow}
\text{(ii)}\handwritten{\Leftrightarrow}\text{(iii)}
\overset{\inserted{\handwritten{\Leftrightarrow\text{(iii bis)}\Leftrightarrow}}}{\mathstrut}
\text{(iv)}\handwritten{\Leftrightarrow}\text{(vi)}\handwritten{\ ;}$}
\tag*{\handwritten{(3)}}
\]
si \deleted{\unsure{$G$|$V$}}\handwritten{$V$} est \deleted{lisse sur $k$ et $W$ géométriquement} \handwritten{\inserted{génériquement}} réduit, on a (iii \deleted{bis}\handwritten{ter})$\handwritten{\Leftrightarrow}\text{(\handwritten{i})}$\handwritten{,} si \deleted{de plus}\handwritten{\inserted{enfin}} $V$ est \deleted{géométriquement} normal \deleted{\illegible{1}}\handwritten{\inserted{on a}} toutes les conditions envisagées sont équivalentes.

On sait déjà par le lemme 1.1. que l'on a $\text{(i)}\handwritten{\Rightarrow}\text{(iv)}\handwritten{\Leftrightarrow}\text{(v\handwritten{i})}$, d'autre part (vi) implique évidemment par définition $\handwritten{N}=\handwritten{M}_a$ et (iv) implique \deleted{(iii bis)} que $\handwritten{\psi}$ est génériquement non ramifié (NB il résulte des hypothèses faites que $X$ est irréductible\deleted{\unsure{s|,}} comme image de $G\mathrm{x}W$ qui l'est), donc $\text{(iv)}\handwritten{\Leftrightarrow}\text{(vi)}\handwritten{\Rightarrow}\text{(iii bis)}$, évidemment $\text{(iii bis)}\handwritten{\Rightarrow}\text{(iii)}$\unsure{,} \deleted{enfin} De plus $\text{(iii)}\handwritten{\Rightarrow}\text{(ii)}$, car comme $\dim X\handwritten{\geq}\dim V$ (1.2.) il s'ensuit que si $\handwritten{\psi}$ est génériquement quasi-fini, \deleted{\unsure{on a}}\handwritten{donc} $\dim X\handwritten{\leq}\dim V$\handwritten{,} \deleted{\unsure{donc}}\handwritten{on a} $\dim X=\dim V$. Enfin, en vertu de 1.2.\handwritten{,} $\text{(ii)}\handwritten{\Rightarrow}\text{(i)}$, ce qui prouve l'équivalence des conditions (i) (ii) (iii) (iii bis) (iv) (vi). Evidemment $\text{(iii ter)}\handwritten{\Rightarrow}\text{(iii bis)}$ et $\text{(iv bis)}\handwritten{\Rightarrow}\text{(iv)}$, donc par ce qui précède
% archive sheet number: 2
